\documentclass[a4paper,11pt]{article}
\usepackage[margin=26mm]{geometry}
\usepackage{amsmath,amssymb,amsthm,booktabs,hyperref}
\newtheorem{theorem}{Theorem}
\newtheorem{lemma}[theorem]{Lemma}
\newtheorem{corollary}[theorem]{Corollary}
\newcommand{\gen}[1]{\langle#1\rangle}
\newcommand{\Sym}{\operatorname{Sym}}
\title{Subgroup intersections under central extensions}
\author{Human Compression project}
\date{Private component for the recovery of L019}
\begin{document}\maketitle
\begin{abstract}
A central extension can remove an intersection defect only when the
corresponding quotient embeds in a quotient of its central kernel. In
particular, an odd factor in the intersection index survives every extension
by central involutions. This supplies a uniform replacement for several
uses of an unrestricted endpoint-untwisting claim in the specified source.
An exact check of its eleven finite drawings verifies ten exclusions and
identifies a positive exception in the remaining drawing. No conclusion
about the source's main stability theorem is claimed here.
\end{abstract}

\section{The difference between two lifts}
All groups in the finite applications are actual permutation groups on the
displayed point set. A marking $s_0,\ldots,s_{r-1}$ is C when, for every
two label sets $I,J$, its generated subgroups satisfy
$G_I\cap G_J=G_{I\cap J}$. The result below does not require involutions
or the string commutation condition.

\begin{lemma}\label{difference}
Let $\pi:\widetilde G\to G$ be a surjective homomorphism with central
kernel $E$. Let $\widetilde U,\widetilde V\leq\widetilde G$, put
$U=\pi(\widetilde U)$, $V=\pi(\widetilde V)$ and
$K_U=\widetilde U\cap E$, $K_V=\widetilde V\cap E$. Then
$P=\pi(\widetilde U\cap\widetilde V)$ is a normal subgroup of
$U\cap V$, and
\[
             (U\cap V)/P\ \hookrightarrow\ E/(K_UK_V).
\]
\end{lemma}
\begin{proof}
For $x\in U\cap V$, choose lifts $u\in\widetilde U$ and
$v\in\widetilde V$ with image $x$. Their difference $uv^{-1}$ lies
in $E$. Replacing the two lifts changes this difference by an element
of $K_UK_V$, so
\[
 \delta(x)=uv^{-1}K_UK_V
\]
is well defined. Centrality of $E$ makes it abelian and makes the displayed
product of subgroups a subgroup. If $u_x,v_x,u_y,v_y$ are the chosen
lifts for $x,y$, the lifts $u_xu_y,v_xv_y$ give
\[
 u_xu_y(v_xv_y)^{-1}
 =u_x(u_yv_y^{-1})v_x^{-1}
 =(u_yv_y^{-1})(u_xv_x^{-1}).
\]
Thus $\delta$ is a homomorphism into $E/(K_UK_V)$.

If $x$ has a common lift in the two subgroups, then $\delta(x)=1$.
Conversely, suppose $uv^{-1}=k_Uk_V$ with $k_U\in K_U$ and
$k_V\in K_V$. The element $k_U^{-1}u=k_Vv$ belongs to both subgroups
and maps to $x$. Hence $\ker\delta=P$. The first isomorphism theorem
gives the asserted normality and embedding.
\end{proof}

\begin{theorem}\label{obstruction}
Suppose a finite marked group $(\widetilde G;\widetilde s_i)$ is C
and $\pi:\widetilde G\to G$ has central kernel $E$. Write
$s_i=\pi(\widetilde s_i)$. For every pair $I,J$ of label sets,
$G_{I\cap J}$ is normal in $G_I\cap G_J$, and
\[
        (G_I\cap G_J)/G_{I\cap J}
        \hookrightarrow E/(E\cap\widetilde G_I)(E\cap\widetilde G_J).
\]
In particular, if $E$ is a $p$-group, then
$[G_I\cap G_J:G_{I\cap J}]$ is a power of $p$.
\end{theorem}
\begin{proof}
Apply Lemma~\ref{difference} to the two marked subgroups upstairs.
Their intersection is $\widetilde G_{I\cap J}$ by the C-property,
and its image is exactly $G_{I\cap J}$ because the homomorphism
sends each of its specified generators to the specified old generator.
The order of the displayed quotient divides the order of the finite
central quotient on the right.
\end{proof}

\begin{corollary}\label{twists}
Let $s_i$ be a marking of a finite group $G$, and let $E$ be a finite
elementary abelian $2$-group. Replace the marks by $(s_i,e_i)$ in
$G\times E$, for arbitrary $e_i\in E$. If an original intersection
index has an odd prime divisor, the resulting marking is not C.
\end{corollary}
\begin{proof}
The subgroup generated by the new marks projects onto $G$. Its kernel
is a subgroup of the central group $E$. The conclusion follows from
Theorem~\ref{obstruction} with $p=2$.
\end{proof}

The corollary permits simultaneous twists at any labels, not merely one
endpoint. For a permutation application, disjoint supports for the original
action and the commuting extra involutions supply precisely the required
direct product. An actual homomorphism to the original group is essential:
an arbitrary quotient or an unproved separation of restriction images
cannot replace it. The theorem does not justify unrestricted untwisting
when the old intersection index is a power of two. In particular it does
not undo the previously verified counterexample to Proposition~2.4 of
\cite{source}.

\section{The finite diagrams in the source}
The source is the exact author edition arXiv:2212.12723v2, 11 January
2023. Its Proposition~2.5, physical page~7, displays eleven rank-four
markings. Number vertices from left to right along the upper row and then
from left to right along the lower row, omitting empty diagram positions.
The attached numeric data record every transposition of every mark under
this convention. Both the primary TeX and the native PDF drawing have
been inspected; no inferred abstract isomorphism identifies the inputs.

Put $U=\gen{s_1,s_2,s_3}$, $V=\gen{s_0,s_1,s_2}$,
$I=U\cap V$ and $M=\gen{s_1,s_2}$. Exact finite calculations give:
\begin{center}
\begin{tabular}{crrrrr}
\toprule
Source diagram&$|U|$&$|V|$&$|I|$&$|M|$&$[I:M]$\\
\midrule
1&4320&10080&240&20&12\\
2&2880&10080&144&24&6\\
3&720&720&120&10&12\\
4&4320&10080&240&20&12\\
5a&10080&10080&480&20&24\\
5b&10080&10080&240&20&12\\
6a&5040&1440&240&20&12\\
6b&5040&1440&120&20&6\\
7&5040&720&144&24&6\\
8&5040&720&144&24&6\\
9&60&720&10&10&1\\
\bottomrule
\end{tabular}
\end{center}
Consequently the first ten markings fail the intersection property, and
so do all their central involution twists covered by
Corollary~\ref{twists}. This is stronger than the required individual
failures. It supplies a valid way to exclude larger markings containing
such a twisted submarking, since every submarking of a C-marking is C.
Each use in a larger permutation graph must still identify its actual
central extension and its specified four labels.

The source's numerical description of diagram~5b overstates $|I|$ by a
factor of two. The exclusion remains valid, with index~12. Diagram~9
requires a different correction: it is a C-marking, not an exclusion.
Its literal permutations are
\[
 s_0=(1\ 2),\quad s_1=(2\ 3)(4\ 5),\quad
 s_2=(3\ 4)(5\ 6),\quad s_3=(4\ 5)(6\ 7).
\]
They generate $S_7$ and have adjacent product orders $6,5,5$.
The two distant pairs involving $s_0$ have disjoint supports; the
remaining distant pair shares the transposition $(4\ 5)$ and has
disjoint remaining transpositions, so the string condition holds.
The finite certificate verifies every one of the 256 subset intersections,
including the empty subset. Thus the claimed universal exclusion in
Proposition~2.5 is false for the displayed diagram~9. Its reported
group $U\cong A_6$ and intersection $I\cong A_5$ have orders 360
and 60; the literal drawing instead gives orders 60 and 10.

\section{Finite verification and remaining scope}
The finite verification uses the following complete procedure. Encode a
permutation by its ordered list of images. Starting with the identity,
repeatedly right-multiply each newly found permutation by each specified
generator and retain every previously unseen product. At termination the
retained set contains the generators, consists of their words, and is
closed under multiplication by them. It is therefore exactly the generated
group. For the table, perform this procedure for $U,V,M$ and intersect
the two literal sets for $U,V$. Their orders never exceed 10080. For
diagram~9, perform it for all sixteen label subsets and test all 256
intersections. The full group has 5040 elements on seven points, proving
that it is $S_7$. This is a finite certificate for the specified inputs,
not a claim about unenumerated degrees or arbitrary markings.

The implementation validates the input matchings, nonidentity and
distinctness of all four marks, the string commutations and every displayed
order. Optimized Python is rejected. A separate GAP calculation checks the
same literal permutation inputs using its subgroup algorithms. The
uniform statement about all central twists follows from the written proof
in the first section, not from sampling a few twists.

The finite exclusions and their use as local submarkings come from
\cite[Proposition~2.5]{source}. The difference homomorphism provides the
present repair of the untwisting interface; no historical priority claim is
made for this elementary argument. The corrected entries are source-audit
findings, not answers to open conjectures. The high-rank existence of
suitable splits, admissible rank extensions, inverse contraction, and the
global counting bijection of L019 remain separate proof obligations.
In particular, uses of the excluded diagram~9 in the source's subsequent
attachment argument must be checked again. No complete original recovery
or refutation of the main stability theorem is asserted.

\begin{thebibliography}{9}
\bibitem{source} P. J. Cameron, M. E. Fernandes and D. Leemans,
\emph{The number of string C-groups of high rank},
arXiv:2212.12723v2, 11 January 2023, Proposition~2.5 and physical page~7.
\end{thebibliography}
\end{document}
